% 数理物理1 数理物理1-8 問題4（061-ring3。当時は 数理物理1-7 問題8・shou07-l2-ring3） 解いた日 2026-10-03
%% kotae
\toi{1}{◎}{}
\begin{align*}
  V&=\frac{\kappa}{2}(x_1^{2}+x_2^{2}+x_3^{2})\\
  &\qquad+\frac{k}{2}\bigl(2x_1^{2}+2x_2^{2}+2x_3^{2}\\
  &\qquad\kakimachigai{-x_1x_2-x_2x_3-x_3x_1}{-2(\cdots)}\bigr)
\end{align*}
\[
  S=\begin{pmatrix}\kappa+2k&-k&-k\\-k&\kappa+2k&-k\\-k&-k&\kappa+2k\end{pmatrix}
\]
\toi{2}{△}{計算}
（ページの上に固有方程式のメモ）
\[ (\kappa+2k-\lambda)^{3}-2k^{3}-3k^{2}(\kappa+2k+\lambda) \]
\[
  U=\frac{1}{\sqrt{3}}\begin{pmatrix}\ee^{\ii\frac{2\pi}{3}}&\ee^{\ii\frac{4\pi}{3}}&\ee^{\ii\frac{6\pi}{3}}\\[2pt]\ee^{\ii\frac{4\pi}{3}}&\ee^{\ii\frac{8\pi}{3}}&\ee^{\ii\frac{12\pi}{3}}\\[2pt]\ee^{\ii\frac{6\pi}{3}}&\ee^{\ii\frac{12\pi}{3}}&\ee^{\ii\frac{18\pi}{3}}\end{pmatrix}
\]
\begin{align*}
  U^{\dagger}SU&=\frac{1}{3}\begin{pmatrix}\zeta^{-1}&\zeta^{-2}&\zeta^{-3}\\\zeta^{-2}&\zeta^{-4}&\zeta^{-6}\\\zeta^{-3}&\zeta^{-6}&\zeta^{-9}\end{pmatrix}\\
  &\qquad\times\begin{pmatrix}\kappa+2k&-k&-k\\-k&\kappa+2k&-k\\-k&-k&\kappa+2k\end{pmatrix}\\
  &\qquad\times\begin{pmatrix}\zeta&\zeta^{2}&\zeta^{3}\\\zeta^{2}&\zeta^{4}&\zeta^{6}\\\zeta^{3}&\zeta^{6}&\zeta^{9}\end{pmatrix}\\
  &=\frac{1}{3}\begin{pmatrix}\zeta^{\fumei}(\kappa+2k)&\cdots\end{pmatrix}\nuke{\lambda_1,\lambda_2,\lambda_3}
\end{align*}
\naze{立式は正しいが，対角成分を出す前に止まっている。$S$ の各行は対角に $\kappa+2k$，両隣に $-k$ なので，$U_{n\pm1,\mu}=\zeta^{\pm\mu}U_{n\mu}$ を使うと $(SU)_{n\mu}=[\kappa+2k-k(\zeta^{\mu}+\zeta^{-\mu})]U_{n\mu}$ とまとまり，$\lambda_{\mu}=\kappa+2k-2k\cos(2\pi\mu/3)$，すなわち $\lambda_1=\lambda_2=\kappa+3k$，$\lambda_3=\kappa$ が出る。行列の積を全部書く必要はない。}
\toi{3}{×}{未着手}
\[
  \dv[2]{\hat{x}_1}{t}=-\frac{\kappa+3k}{m}\hat{x}_1,\qquad x=
\]
（図：3個の質点を三角形に並べ，ばねでつないだ図）
\naze{1つの式だけで止まっている。$\partial V/\partial x_n=(S\vect{x})_n$ より $m\,\dd^{2}\vect{x}/\dd t^{2}=-S\vect{x}$，$UU^{\dagger}=I$ を挟んで左から $U^{\dagger}$ をかけると $m\,\dd^{2}\hat{x}_{\mu}/\dd t^{2}=-\lambda_{\mu}\hat{x}_{\mu}$ となり，各モードが独立に単振動する。$\omega_1=\omega_2=\sqrt{(\kappa+3k)/m}$，$\omega_3=\sqrt{\kappa/m}$。}
\toi{4}{×}{未着手}
\mikai
%% kizuki
\item 環状の系の $U^{\dagger}SU$ は，行列の積を書き下さずに「$S$ の第 $n$ 行は対角 $\kappa+2k$，両隣 $-k$」と「$U_{n\pm1,\mu}=\zeta^{\pm\mu}U_{n\mu}$」から $(SU)_{n\mu}=\lambda_{\mu}U_{n\mu}$ を出す。$U^{\dagger}U=I$ と合わせれば $U^{\dagger}SU$ は対角で，対角成分が $\lambda_{\mu}$ になる。固有方程式の3次式を解くのは遠回りである。
\item 基準座標への書き換えは「$m\ddot{\vect{x}}=-S\vect{x}$ に左から $U^{\dagger}$ をかけ，$\vect{x}=U\hat{\vect{x}}$ を入れる」の1行で，全モードの式と $\omega_{\mu}=\sqrt{\lambda_{\mu}/m}$ が一度に出る。1つのモードだけ書いて終えない。
